\chapter{優先度に基づくMCL間のパーティクル共有と自己位置推定の復帰}
\label{cha:simulation_experiments}
\graphicspath{{figures/06}}
\section{はじめに}

\section{提案手法の概要}

第7章では，LiDAR，地磁気およびWi-Fiの観測情報を用いて各MCLの推定状態を評価し，位置推定に使用するMCLを切り替える手法について述べた．しかし，異常と判定されたMCLが誤った位置に収束している場合，位置推定に使用するMCLを切り替えるだけでは，異常となったMCL自身の推定状態を回復させることはできない．

そこで本研究では，LiDAR MCL，地磁気MCL，Wi-Fi MCLおよびこれらの観測情報を統合したMulti MCLを並列に動作させ，各MCLの推定状態に応じて粒子を共有する手法を提案する．各MCLの観測整合性および粒子分布の状態を優先度として表し，その優先度に基づいて外部粒子の導入割合と取り込み元を決定する．これにより，推定状態が悪化したMCLに対して，他のMCLが保持する位置仮説を導入し，自己位置推定の復帰を図る．

LiDAR MCLおよびMulti MCLでは，自身の粒子群による位置推定を基本とし，推定状態の悪化に応じて他のMCLの粒子を導入する．一方，地磁気MCLおよびWi-Fi MCLでは，観測情報の識別性が十分でない場合を考慮し，通常状態でもリサンプリング時に外部粒子を導入する．このように，各MCLの特性に応じて粒子共有の方針を変えることで，正常時の推定状態を維持しながら，異常状態からの復帰を図る．

各MCLの優先度を $p_i$ とする．優先度は数値が小さいほど位置推定への採用を優先し，異常状態を $p_i=100$ とする．本手法では，LiDAR MCLおよびMulti MCLの基準優先度を $0$，地磁気MCLを $80$，Wi-Fi MCLを $90$ とする．これは，推定状態が良好なLiDAR MCLおよびMulti MCLを優先し，地磁気MCLおよびWi-Fi MCLを補助的な位置仮説の供給源として利用するためである．


\section{推定状態に応じた優先度制御}

\subsection{成功・失敗の確率分布に基づく失敗指標の算出}

本研究では，LiDAR MCLおよびMulti MCLの推定状態を評価するため，第7章で述べた平滑化後の平均LiDAR尤度および粒子分布の広がりを使用する．第7章では，これらの指標をあらかじめ設定した閾値と比較して異常状態を判定した．これに対し，本章では，各指標を成功時および失敗時の確率分布に基づく失敗指標へ変換し，推定状態の悪化の程度を評価する．

失敗指標の算出には，事前に取得した走行データを使用する．LiDAR MCLおよびMulti MCLのそれぞれについて，推定位置と基準位置との二次元距離が $1\,\mathrm{m}$ 未満の場合を成功状態 $S$，$1\,\mathrm{m}$ 以上の場合を失敗状態 $F$ とする．

各MCLについて，成功状態および失敗状態における平均LiDAR尤度と粒子分布の広がりを抽出し，それぞれヒストグラムを作成する．各ヒストグラムでは，観測値を一定幅の区間に分割し，成功時と失敗時の出現頻度から条件付き確率分布を推定する．

指標の値を $x$ としたとき，失敗指標 $q(x)$ を次式で定義する．

\[
q(x)=
\frac{p(x\mid F)}
{p(x\mid S)+p(x\mid F)}
\]

ここで，$p(x\mid S)$ および $p(x\mid F)$ は，それぞれ成功状態および失敗状態における指標 $x$ の条件付き確率密度である．本研究では，成功状態と失敗状態の事前確率をそれぞれ $0.5$ とする．したがって，$q(x)$ は等事前確率を仮定した場合の失敗状態の事後確率に相当する．

ヒストグラムの区間を $B_k$ とし，区間 $B_k$ に含まれる成功データ数を $n_{S,k}$，失敗データ数を $n_{F,k}$ とする．また，成功データの総数を $N_S$，失敗データの総数を $N_F$ とすると，各区間における失敗指標は次式で算出できる．

\[
q_k=
\frac{n_{F,k}/N_F}
{n_{S,k}/N_S+n_{F,k}/N_F}
\]

このように，成功時と失敗時のデータ数をそれぞれの総数で正規化することで，走行データに含まれる成功データと失敗データの件数が異なる場合にも，等事前確率に基づいて失敗指標を算出する．

ただし，ヒストグラムから直接算出した失敗指標は，各区間に含まれるデータ数の違いによって急激に変動する場合がある．そこで本研究では，失敗指標に対してガウシアン平滑化を適用する．

区間 $B_k$ における平滑化後の失敗指標を $\tilde{q}_k$ とすると，次式で表される．

\[
\tilde{q}_k=
\frac{
\displaystyle\sum_j G_{\sigma}(k-j)m_jq_j
}{
\displaystyle\sum_j G_{\sigma}(k-j)m_j
}
\]

ここで，$G_{\sigma}$ はガウシアンカーネル，$\sigma$ は平滑化幅，$m_j$ は失敗指標が定義されている区間で $1$，未定義の区間で $0$ となる値である．これにより，未観測区間を除外しながら，近傍区間の失敗指標を用いて平滑化を行う．なお，成功時および失敗時の出現頻度や条件付き確率分布には平滑化を適用せず，失敗指標のみを平滑化する．

以上の処理により，LiDAR MCLおよびMulti MCLのそれぞれについて，平均LiDAR尤度および粒子分布の広がりに対応する参照テーブルを作成する．参照テーブルには各区間の下限値，上限値および平滑化後の失敗指標を格納する．

実行時には，各MCLで得られた指標に対応する区間を参照し，平滑化後の失敗指標を取得する．平均LiDAR尤度に対応する失敗指標を $q_{\ell,i,t}$，粒子分布の広がりに対応する失敗指標を $q_{s,i,t}$ とする．

\subsection{失敗指標に基づく優先度の決定}

前項で算出した失敗指標を用いて，LiDAR MCLおよびMulti MCLの優先度を決定する．

平均LiDAR尤度および粒子分布の広がりに対応する失敗指標の平均を，次式で定義する．

\[
\bar{q}_{i,t}=
\frac{q_{\ell,i,t}+q_{s,i,t}}{2}
\]

ここで，$\bar{q}_{i,t}$ はMCL $i$ の総合的な失敗指標である．この値が大きいほど，平均LiDAR尤度または粒子分布の広がりから判断される推定状態が悪いことを表す．

本研究では，総合的な失敗指標に応じて優先度を変化させる．通常状態における優先度は，次式で求める．

\[
p_{i,t}=
\operatorname{clip}
\left(
\operatorname{round}
\left[
b_i+(u-b_i)
\frac{\bar{q}_{i,t}}{\tau_q}
\right],
b_i,u
\right)
\]

ここで，$b_i$ は各MCLの基準優先度，$u=80$ は通常状態における優先度の上限，$\tau_q=0.8$ は失敗指標の判定閾値である．

この式では，失敗指標が小さい場合には優先度が基準値に近くなり，失敗指標の増加に伴って優先度の数値が大きくなる．これにより，推定状態の悪化に応じて，他のMCLから導入する粒子の割合を段階的に増加させることができる．

ただし，いずれかの失敗指標が閾値 $\tau_q$ 以上となった場合には，LiDAR MCLとMulti MCLで異なる処理を行う．LiDAR MCLでは，推定状態が異常であると判断し，優先度を $100$ に設定する．一方，Multi MCLでは，LiDAR以外の観測情報も利用していることを考慮し，この条件のみでは異常状態とせず，優先度を通常状態の上限である $80$ とする．

また，失敗指標の算出に必要な情報が無効な場合には，両MCLとも優先度を $100$ に設定する．

\section{各MCLの優先度制御}

LiDAR MCLおよびMulti MCLでは，前項で求めた失敗指標に加え，第7章で説明したWi-Fi観測との整合性および他のMCLとの位置整合性を用いて優先度を決定する．

Wi-Fi観測との整合性は，各MCLの粒子群に対するWi-Fi信号強度の平均尤度に基づいて判定する．Wi-Fi観測の更新時に平均尤度を閾値と比較し，次の更新までは判定結果を保持する．平均尤度が閾値未満となった場合，または初回の判定が完了していない場合には，異常状態とする．

また，位置整合性については，LiDAR MCLまたはMulti MCLの推定位置が，Wi-Fi MCLおよび地磁気MCLの95%誤差楕円内に含まれているかを評価する．両MCLが位置整合性の評価に利用可能であり，自身の推定位置が両方の誤差楕円外となった場合には，異常状態とする．

これらの判定のいずれかが異常となった場合，LiDAR MCLおよびMulti MCLの優先度を $100$ に設定する．

一方，地磁気MCLおよびWi-Fi MCLについては，失敗指標に基づく連続的な優先度制御を行わず，観測整合性に基づいて優先度を決定する．地磁気MCLでは，平均地磁気尤度が閾値以下となった場合，またはWi-Fi平均尤度による異常判定が成立した場合に優先度を $100$ とし，それ以外では基準優先度の $80$ を使用する．Wi-Fi MCLでは，Wi-Fi平均尤度が閾値未満となった場合に優先度を $100$ とし，それ以外では基準優先度の $90$ を使用する．

以上のように，各MCLの観測特性に応じて優先度を決定することで，位置推定に利用するMCLの選択に加え，粒子共有における外部粒子の導入割合および取り込み元の決定に優先度を利用する．

\section{優先度に基づくMCL間の粒子共有}

\subsection{外部粒子の導入割合の決定}

本研究では，各MCLの優先度に基づいて，他のMCLから導入する粒子の割合を決定する．外部粒子の導入割合を $\rho_i$ とし，次式で求める．

\[
\rho_i=
\begin{cases}
0.2,
& p_i=100,\\[3pt]
0.2\,
\operatorname{clip}
\left(
\dfrac{p_i-b_i}{80-b_i},0,1
\right),
& i\in\{\mathrm{LiDAR},\mathrm{Multi}\},
\quad p_i\ne100,\\[8pt]
0.2,
& i\in\{\mathrm{Mag},\mathrm{WiFi}\},
\quad p_i\ne100.
\end{cases}
\]

LiDAR MCLおよびMulti MCLでは，優先度が基準値である場合に外部粒子を導入せず，自身の粒子のみを使用する．一方，推定状態の悪化に伴って優先度の数値が増加すると，外部粒子の導入割合を増加させる．これにより，自身の位置仮説を維持しながら，他のMCLが保持する位置仮説を取り込むことができる．

地磁気MCLおよびWi-Fi MCLでは，通常状態でも外部粒子を $20\%$ 導入する．地磁気およびWi-Fiの観測情報は，異なる位置で類似した観測値が得られる場合があるため，単独観測では位置を一意に特定することが困難となる可能性がある．そこで，他のMCLが保持する位置仮説を継続的に取り込むことで，位置推定の曖昧さを補う．

なお，いずれのMCLにおいても，異常状態における外部粒子の導入割合は $20\%$ とし，異常状態への移行によって導入割合をさらに増加させる設定にはしていない．

\subsection{取り込み元MCLの選択と粒子数の配分}

MCL $i$ が保持する粒子数を $N_i$ としたとき，外部から導入する粒子数 $K_i$ は次式で求める．

\[
K_i=\operatorname{round}(N_i\rho_i)
\]

本手法では，自身の粒子から $N_i-K_i$ 個，他のMCLから合計 $K_i$ 個を抽出することで，共有後の総粒子数を一定に維持する．

外部粒子の取り込み元は，自身以外のMCLのうち，優先度が $0$ 以上 $95$ 以下であり，粒子情報が有効なものとする．利用可能な取り込み元の集合を $\mathcal{D}_i$ とし，MCL $j$ からの粒子の採用比率 $a_{ij}$ を次式で定義する．

\[
a_{ij}=
\frac{100-p_j}
{\displaystyle\sum_{k\in\mathcal{D}_i}(100-p_k)}
\]

この式により，優先度の数値が小さいMCLほど採用比率が大きくなる．したがって，良好な推定状態にあるMCLの位置仮説を優先的に導入できる．

各取り込み元からの粒子数は，外部粒子数 $K_i$ と採用比率 $a_{ij}$ に基づいて決定する．ただし，粒子数は整数であるため，算出した採用数の整数化には最大剰余法を用いる．これにより，各取り込み元への割当粒子数の合計を $K_i$ に一致させる．

なお，利用可能な取り込み元が存在しない場合には，外部粒子を導入せず，自身の粒子のみを使用する．

\subsection{粒子共有を伴うリサンプリング}

本手法では，外部粒子の導入をリサンプリング処理と組み合わせて実行する．

各取り込み元からの粒子の抽出には，粒子重みに基づく系統的リサンプリングを使用する．これにより，取り込み元MCLにおいて観測との整合性が高い粒子ほど選択されやすくなる．

粒子共有は，MCL間での粒子数の配分と，各MCL内での粒子抽出という二段階の処理で構成される．まず，各MCLの優先度を用いて取り込み元と導入粒子数を決定する．次に，各取り込み元の粒子重みに基づいて指定個数の粒子を抽出し，自身の粒子と組み合わせることで粒子集合を再構成する．

また，粒子共有を伴うリサンプリングは，各MCLの推定状態および観測更新に応じて実行する．

LiDAR MCL，Multi MCLおよび地磁気MCLでは，通常状態において，有効サンプルサイズの低下に加え，前回のリサンプリングからの移動距離または回転角が所定の条件を満たした場合にリサンプリングを実行する．一方，異常状態では，これらの条件によらずリサンプリングを実行する．

Wi-Fi MCLでは，Wi-Fi観測の更新時に，有効サンプルサイズの低下または異常状態を条件としてリサンプリングを実行する．これにより，Wi-Fi観測の更新周期に合わせて粒子集合を再構成する．

以上の処理により，LiDAR MCLおよびMulti MCLでは，正常時の粒子分布を維持しながら，推定状態が悪化した場合に他のMCLの位置仮説を導入する．また，地磁気MCLおよびWi-Fi MCLでは，リサンプリング時に継続的な粒子共有を行うことで，単独観測による位置推定の曖昧さを補う．

このように，各MCLの推定状態を反映した優先度に基づいて粒子を共有することで，誤った位置に収束したMCLに対して他のMCLが保持する位置仮説を導入し，自己位置推定の復帰を図る．
\section{自己位置推定の実験}
\subsection{実験方法}

\subsection{実験結果}

\section{考察}



\section{おわりに}





\section{シミュレーション環境の構築}
\label{sec:sim_environment}

本研究では，LiDAR，地磁気，Wi-Fi，およびこれらを組み合わせた Multi MCL を並列に動作させ，各 MCL の推定状態に応じて粒子を共有する手法を提案する．各 MCL の観測整合性および粒子分布の状態を優先度として表し，その優先度に基づいて外部粒子の導入量と取り込み元を決定する．

LiDAR MCL および Multi MCL では，自身の粒子フィルタによる推定を基本とし，推定状態の悪化に応じて他の MCL の位置仮説を導入する．一方，地磁気 MCL および Wi-Fi MCL は，単独では位置推定の精度や識別性が十分でない場合を想定し，通常状態でもリサンプリングのたびに外部粒子を導入する．このように，各 MCL の特性に応じて粒子共有の方針を変えることで，推定の継続性と，推定状態が悪化した際の回復を図る．

各 MCL の優先度を $p_i$ とする．優先度は数値が小さいほど採用を優先し，異常状態を $p_i=100$ で表す．基準優先度は，LiDAR MCL と Multi MCL を $0$，地磁気 MCL を $80$，Wi-Fi MCL を $90$ とする．これは，良好な状態にある LiDAR・Multi MCL の位置仮説を優先的に利用し，地磁気・Wi-Fi MCL を補助的な位置仮説の供給源として扱う設計である．

LiDAR MCL および Multi MCL の優先度は，平均 LiDAR 尤度と粒子分布の広がりから求める．平均 LiDAR 尤度には一次ローパスフィルタを適用し，一時的な観測変動の影響を抑える．
\begin{equation}
  \bar{\ell}_{i,t}
  =
  (1-\alpha)\bar{\ell}_{i,t-1}
  +\alpha\ell_{i,t}.
  \label{eq:lidar_likelihood_lpf}
\end{equation}
ここで，$\ell_{i,t}$ は時刻 $t$ における平均 LiDAR 尤度，$\bar{\ell}_{i,t}$ は平滑化後の値である．

粒子分布の広がりには，粒子の二次元位置から等重みで算出した共分散行列 $\Sigma_{i,t}$ を用い，次式で表される誤差楕円の長半径を使用する．
\begin{equation}
  s_{i,t}
  =
  \sqrt{\chi^2_{2,0.95}\lambda_{\max}(\Sigma_{i,t})}.
  \label{eq:particle_spread}
\end{equation}
ここで，$\lambda_{\max}$ は最大固有値であり，$\chi^2_{2,0.95}=5.991$ とする．この値は粒子分布の広がりを評価するための尺度として使用する．

平滑化後の LiDAR 尤度と分布の広がりを，それぞれ事前に用意した参照テーブルにより，失敗指標 $q_{\ell,i,t}$，$q_{s,i,t}$ に変換する．参照テーブルには各指標の区間に対応する平滑化済みの失敗指標が格納されており，LiDAR 尤度の低下や粒子分布の拡大に伴う推定状態の悪化を評価する．各指標の上下限を超えた場合には，良好側を $0$，悪化側を $1$ とする．

両指標の平均を
\begin{equation}
  \bar{q}_{i,t}
  =
  \frac{q_{\ell,i,t}+q_{s,i,t}}{2}
  \label{eq:mean_failure_indicator}
\end{equation}
とし，通常時の優先度を
\begin{equation}
  p_{i,t}
  =
  \operatorname{clip}
  \left(
    \operatorname{round}
    \left[
      b_i+(u-b_i)\frac{\bar{q}_{i,t}}{\tau_q}
    \right],
    b_i,u
  \right)
  \label{eq:normal_priority}
\end{equation}
で算出する．ここで，$b_i$ は基準優先度，$u=80$ は通常時の上限，$\tau_q=0.8$ は失敗指標の判定閾値である．

ただし，いずれかの失敗指標が閾値以上となった場合，LiDAR MCL では優先度を $100$ とする．Multi MCL では，LiDAR 以外の観測も使用することを考慮し，この条件だけでは異常状態とせず，優先度を通常時の上限 $80$ とする．指標の算出に必要な情報が無効な場合は，両 MCL とも優先度を $100$ とする．

さらに，Wi-Fi 観測との整合性，および他の MCL との位置整合性によって異常判定を補完する．Wi-Fi 観測との整合性は，各 MCL が保持する粒子群の平均 Wi-Fi 尤度を用いて評価する．Wi-Fi 観測の更新時に閾値判定を行い，次の更新まではその結果を保持する．平均尤度が閾値未満の場合，または初回の判定が行われていない場合は，異常状態とする．

位置整合性については，LiDAR MCL または Multi MCL の推定位置が，Wi-Fi MCL と地磁気 MCL の誤差楕円内にあるかを評価する．取り込み元として利用可能な両 MCL に対し，自身の推定位置が両方の楕円外にある場合は異常状態とする．楕円内外の判定には，共分散行列を正則化したマハラノビス距離を用いる．Wi-Fi 尤度判定または位置整合性判定が異常となった場合，LiDAR MCL と Multi MCL の優先度を $100$ に設定する．

地磁気 MCL では，平均地磁気尤度が閾値以下の場合，または Wi-Fi 尤度判定が異常の場合に優先度を $100$ とし，それ以外では基準優先度 $80$ を用いる．Wi-Fi MCL では，Wi-Fi 尤度判定に応じて基準優先度 $90$ と異常時の $100$ を切り替える．

粒子共有では，各 MCL の優先度に応じて，粒子集合の一部を他の MCL 由来の粒子で置き換える．外部粒子の導入割合を $\rho_i$ とし，次式で決定する．
\begin{equation}
  \rho_i
  =
  \begin{cases}
    0.2,
      & p_i=100, \\[3pt]
    0.2\,
    \operatorname{clip}
    \left(
      \dfrac{p_i-b_i}{80-b_i},0,1
    \right),
      & i\in\{\mathrm{LiDAR},\mathrm{Multi}\},
        \quad p_i\ne100, \\[8pt]
    0.2,
      & i\in\{\mathrm{Mag},\mathrm{WiFi}\},
        \quad p_i\ne100.
  \end{cases}
  \label{eq:particle_sharing_ratio}
\end{equation}

LiDAR MCL および Multi MCL は，優先度が基準値にある場合には自身の粒子のみを使用する．推定状態が悪化して優先度の数値が増加すると，外部粒子の割合を最大 $20\%$ まで増やす．これにより，自身の推定を維持しながら，状態の悪化に応じて他の MCL の位置仮説を取り込む．

地磁気 MCL および Wi-Fi MCL は，通常状態でも外部粒子を $20\%$ 導入する．これは，単独観測による位置推定の曖昧さを，他の MCL が保持する位置仮説によって補うためである．なお，いずれの MCL も異常時の共有割合は $20\%$ であり，異常状態への移行によって共有割合をさらに増やす設定にはしていない．

MCL $i$ の粒子数を $N_i$ とすると，外部から導入する粒子数は
\begin{equation}
  K_i
  =
  \operatorname{round}(N_i\rho_i)
  \label{eq:external_particle_count}
\end{equation}
となる．自身から $N_i-K_i$ 個，他の MCL から合計 $K_i$ 個を抽出することで，総粒子数を維持する．

取り込み元は，自身以外の MCL のうち，優先度が $0$ 以上 $95$ 以下であり，粒子情報が有効なものとする．利用可能な取り込み元の集合を $\mathcal{D}_i$ とし，MCL $j$ からの採用比率を
\begin{equation}
  a_{ij}
  =
  \frac{100-p_j}
       {\displaystyle\sum_{k\in\mathcal{D}_i}(100-p_k)}
  \label{eq:donor_allocation_ratio}
\end{equation}
で決定する．これにより，優先度の数値が小さい MCL ほど多くの粒子を提供する．採用数の整数化には最大剰余法を用い，外部粒子数の合計を $K_i$ に一致させる．利用可能な取り込み元が存在しない場合は，自身の粒子のみを使用する．

各取り込み元の内部では，粒子重みに基づく系統的リサンプリングによって指定個数の粒子を抽出する．したがって，本手法は，MCL 間では優先度に基づいて採用数を配分し，各 MCL 内では観測に基づく粒子重みによって位置仮説を選択する，二段階の粒子選択として構成される．

共有を伴うリサンプリングは，有効サンプルサイズに基づいて実行する．LiDAR・地磁気・Multi MCL では，通常時は有効サンプルサイズの低下に加え，前回のリサンプリングからの移動距離または回転角を実行条件とする．異常状態ではこれらの条件によらず実行する．Wi-Fi MCL では，Wi-Fi 観測の更新時に，有効サンプルサイズの低下または異常状態を条件として実行する．これにより，各 MCL の観測更新と推定状態に応じたタイミングで粒子集合を再構成する．



\section{考察}
